\begin{theorem}
	Пусть расширение $L \supset F$ сепарабельно. Тогда следующие условия эквивалентны:
	\begin{enumerate}
		\item Расширение $L$ нормально
		\item $L$ является полем разложения некоторого многочлена $f \in F[x]$
		\item $|Aut_F{L}| = [L : F]$
		\item $L^{Aut_F{L}} = F$
	\end{enumerate}
\end{theorem}

\begin{proof}~
	\begin{itemize}
		\item[$1 \Rightarrow 2$] Рассмотрим $\gamma \in L$ такой, что $L = F(\gamma)$. Все сопряженные с $\gamma$ элементы лежат в $L$, поэтому минимальный многочлен $m_\gamma \in F[x]$ раскладывается на линейные сомножители над $L$. Поскольку $L = F(\gamma)$, то это минимальное поле, удовлетворяющее этому условию. Значит, $L$ "--- это поле разложения многочлена $m_\gamma$.
		\item[$2 \Rightarrow 1$]Пусть $L$ "--- поле разложения многочлена $f \in F[x]$. Рассмотрим произвольный элемент $\gamma \in L$ с минимальным многочленом $m_\gamma \in F[x]$. Пусть $\delta \in L$ "--- сопряженный с ним элемент. Тогда $F(\gamma) \cong F[x]/(m_\gamma) \cong F(\delta)$, причем изоморфизм $\varphi: F(\gamma) \to F(\delta)$ сохраняет $F$. Пока оставим без доказательства, что тогда $\varphi$ можно продолжить до гомоморфизма $\widetilde{\varphi}: L \to \overline{F}$. Но тогда $\widetilde{\varphi}$ переводит корни многочлена $f$ в сопряженные с ними элементы, поэтому $\widetilde{\varphi}(L) \subset L$ и $\delta = \varphi(\gamma) = \widetilde{\varphi}(\gamma) \in L$.
		\item[$1 \Rightarrow 3$] Мы знаем, что $|Aut_{F}L| = |\{\alpha \in L : \alpha\text{ сопряжен с }\gamma\}|$. Тогда, поскольку расширение $L$ нормально, $|Aut_{F}{L}| = [L : F]$.
		\item[$3 \Rightarrow 4$] Пусть $L^{Aut_FL} = K \supset F$, тогда $Aut_FL = Aut_KL$. Мы уже знаем, что всегда выполнено $|Aut_KL| \le [L : K]$. Но тогда $[L : K] \ge |Aut_KL| = |Aut_FL| = [L : F]$, поэтому $[L : K] = [L : F]$ и $K = F$.
		\item[$4 \Rightarrow 1$] Зафиксируем произвольный элемент $\gamma \in L$. Рассмотрим следующий многочлен: $f(x) := \prod_{g \in Aut_FL}(x - g(\gamma)) \in L[x]$. Для любого автоморфизма $h \in Aut_FL$ выполнено $h(f)(x) = \prod_{g \in Aut_FL}(x - h(g(\gamma))) = \prod_{\widetilde{g} \in Aut_FL}(x - \widetilde{g}(\gamma)) = f(x)$. Значит, все коэффициенты многочлена $f$ лежат в $L^{Aut_FL} = F$. Тогда, по построению, любой сопряженный с $\gamma$ элемент является корнем $f$, поэтому он имеет вид $g(\gamma)$ для некоторого $g \in Aut_FL$ и лежит в $L$.\qedhere
	\end{itemize}
\end{proof}

\begin{definition}
	Конечное нормальное сепарабельное расширение $L \supset F$ поля $F$ называется \textit{расширением Галуа}.
\end{definition}

\begin{lemmanote}
	Аналогично доказанному выше, для необязательно сепарабельного расширения $L \supset F$ эквиваленты следующие условия:
	\begin{enumerate}
		\item Расширение $L$ нормально
		\item $L$ является полем разложения некоторого многочлена $f \in F[x]$
		\item Любой гомоморфизм полей $\varphi: L \to \overline{F}$, сохраняющий $F$, является автоморфизмом поля $L$
	\end{enumerate}
	
	Приведем доказательство эквивалентности:
	\begin{itemize}
		\item[$1 \Rightarrow 2$] Пусть элементы $\gamma_1, \dots, \gamma_n \in L$ образуют базис в $L$ над $F$. Рассмотрим их минимальные многочлены $m_{\gamma_1}, \dotsc, m_{\gamma_n} \in F[x]$, тогда, в силу нормальности, все корни этих многочленов лежат в $L$, причем $L$ "--- минимальное по включению поле, содержащее их. Значит, $L$ "--- поле разложения многочлена $m_{\gamma_1}\dotsm m_{\gamma_n}$.
		
		\item[$2 \Rightarrow 3$] Пусть гомоморфизм $\varphi: L \to \overline{F}$ сохраняет $F$. Тогда $\varphi$ осуществляет перестановку корней многочлена $f$, полем разложения которого является $L$, поэтому $\varphi(L) = L$ и $\varphi$ "--- автоморфизм.
		
		\item[$3 \Rightarrow 1$] Пусть $\gamma \in L$ "--- произвольный элемент, $\widetilde\gamma \in \overline{F}$ "--- сопряженный с $\gamma$ элемент. Тогда существует изоморфизм $\varphi: F(\gamma) \to F(\widetilde\gamma)$, сохраняющий $F$. Такой изоморфизм можно продолжить до гомоморфизма $\widetilde\varphi : L \to \overline{F}$, который является автоморфизмом поля $L$ по условию. Значит, $\widetilde\gamma \in L$.
	\end{itemize}

	Действительно, если $L$ "--- поле разложения многочлена $f \in F[x]$, то вложение $\varphi: L \to \overline{F}$, сохраняющее $F$, осуществляет перестановку корней многочлена $f$, поэтому $\varphi(L) = L$ и $\varphi$ является автоморфизмом поля $L$.
\end{lemmanote}